\subsection{W-19 (TWA)}
Bezpečnost webových aplikací. Rizika a jejich opatření. Autentizace a autorizace uživatele.

\subsubsection*{Autentizace a její dělení}

\begin{itemize}
  \item Něco ví - jméno a heslo, PIN, ověřovací otázka
  \item Něco má - telefon, token, certifikát
  \item Někým je - otisk prstu, snímek duhovky, DNA
\end{itemize}

HTTP Basic:
\begin{itemize}
  \item uživatelské jméno a heslo se posílá v base64 
  \item není šifrované - je vždy nutné použít HTTPS
  \item Authorization: Basic base64(username:password)
\end{itemize}

HTTP Digest: 
\begin{itemize}
  \item Celý protokol - klient pošle požadavek že se chce přihlásit, server odpoví 401 s "nonce" - náhodný token (a další parametry)
  \item Klient spočítá hash MD5 z nonce, uživatelského jména a hesla a url a odešle ho na server
  \item Server si pomocí hashe ověří, že klient zná co potřebuje
  \item Stejně jako HTTP Basic je to zastaralá metoda
\end{itemize}


API klíč:
\begin{itemize}
  \item Slouží k autentizaci služeb/aplikací.
  \item Aplikace klientovi vydá náhodný token, který poté aplikace použije v hlavičce (případně jinde)
  \item Token musí být náhodný, neuhodnutelný a \textbf{odvolatelný} - lze ho zrušit.
\end{itemize}

OAuth:
\begin{itemize}
  \item Protokol pro autentizaci pomocí třetí strany
\end{itemize}

\subsubsection*{Autorizace}
\begin{itemize}
  \item Role - .hasRole, staticky definované role
  \item ACL - Access Control List - práva jsou přiřazována jednotlivým uživatelům pro konkrétní objekty 
  \item RBAC - Role Based Access Control - uživatelé mají role a role mají oprávnění
\end{itemize}

\subsubsection*{Zranitelosti}

\begin{itemize}
  \item SQL Injection
  \item XSS - Cross Site Scripting
  \item Remote code execution 
  \item CSRF - Cross Site Request Forgery - zneužití uživatelské relace - např. pošlu uživateli obrázek s src="http://mojeaplikace.cz/transfer?amount=1000\&to=mojeid"
\end{itemize}
